
\documentclass[11pt]{article}
\usepackage{amsmath,amssymb,amsthm,mathtools,enumitem,geometry}
\geometry{margin=1in}

\title{Step 73: Weil Positivity versus the Douglas Feature Gate}
\author{Six Birds / Formed-Layer Membrane Program}
\date{}

\newtheorem{theorem}{Theorem}
\newtheorem{lemma}{Lemma}
\newtheorem{proposition}{Proposition}
\newtheorem{definition}{Definition}
\newtheorem{warning}{Warning}
\newtheorem{obligation}{Obligation}

\newcommand{\AZ}{\mathsf A_Z}
\newcommand{\Kcmp}{\mathsf K_{\mathrm{cmp}}}
\newcommand{\Yminus}{Y^-}

\begin{document}
\maketitle

\begin{abstract}
This note compares the framework's Douglas-contractivity obligation
\[
        \AZ \preceq \Kcmp
        \quad\Longleftrightarrow\quad
        V_Z=T W_{\mathrm{cmp}},\quad \|T\|\le1
\]
with classical Weil positivity.  The conclusion is precise: a full-domain
Weil positivity identity for the exact same anti-invariant response space is
equivalent to Loewner domination, hence to Douglas contraction.  However,
classical Weil positivity as usually encountered is a positivity criterion on
a test-function class, often with signed explicit-formula components.  To
become the framework's O1 bridge it must be upgraded by: exact response-space
identification, density/closure of the test class, positive carrier-side feature
decomposition, completion/tail records, and no-smuggling of zero-side data into
the carrier feature map.
\end{abstract}

\section{Objects}

Let \(\Yminus\) be the anti-invariant response space for a formed zeta carrier.
Let
\[
        V_Z:\Yminus\to \mathcal Z,
        \qquad
        \AZ=V_Z^*V_Z
\]
be the zero-side anti-invariant displacement readout.  Let
\[
        W_{\mathrm{cmp}}:\Yminus\to \mathcal E_{\mathrm{cmp}},
        \qquad
        \Kcmp=W_{\mathrm{cmp}}^*W_{\mathrm{cmp}}
\]
be the completed carrier-side feature map.

The framework's explicit-formula domination gate is
\[
        \AZ\preceq \Kcmp.
\]
By the Douglas lemma this is equivalent to the existence of a contraction
\[
        V_Z=T W_{\mathrm{cmp}},
        \qquad \|T\|\le1.
\]

\section{Weil positivity as a quadratic form}

Classical Weil-style criteria are naturally expressed as positivity of a
quadratic form
\[
        Q_W(y)\ge0
        \qquad (y\in\mathcal T),
\]
where \(\mathcal T\) is a declared test class.  In the membrane notation the
desired full anti-invariant form is
\[
        Q_{\mathrm{mem}}(y)
        =
        \langle y,(\Kcmp-\AZ)y\rangle.
\]

\begin{theorem}[Full-domain equivalence]
Suppose \(\Yminus\) is a Hilbert space and \(\AZ,\Kcmp\) are bounded positive
operators on \(\Yminus\).  Then the following are equivalent:
\begin{enumerate}[label=(\roman*)]
\item \( \AZ\preceq \Kcmp\);
\item \( \langle y,(\Kcmp-\AZ)y\rangle\ge0\) for every \(y\in\Yminus\);
\item there exists a contraction \(T\) such that \(V_Z=T W_{\mathrm{cmp}}\),
      for any factorizations \(\AZ=V_Z^*V_Z\) and \(\Kcmp=W_{\mathrm{cmp}}^*W_{\mathrm{cmp}}\).
\end{enumerate}
\end{theorem}

\begin{proof}
The equivalence of (i) and (ii) is the definition of the Loewner order.
The equivalence of (i) and (iii) is Douglas factorization: \(V_Z^*V_Z\preceq
W_{\mathrm{cmp}}^*W_{\mathrm{cmp}}\) iff
\(\operatorname{Ran}(V_Z^*)\subseteq\operatorname{Ran}(W_{\mathrm{cmp}}^*)\)
with the corresponding factor bounded by \(1\), equivalently \(V_Z=T
W_{\mathrm{cmp}}\) for a contraction \(T\).
\end{proof}

\section{The exact gap}

The theorem says that \emph{full-domain} Weil positivity, once identified with
\(Q_{\mathrm{mem}}\), is exactly the Douglas gate.  The actual gap is not formal
positivity; it is earning the identification.

\begin{obligation}[Weil-to-Douglas upgrade]
To use Weil positivity as O1, one must provide:
\begin{enumerate}[label=O1.\arabic*]
\item a precise anti-invariant response space \(\Yminus\);
\item a zero readout \(V_Z\) whose ledger is \(\AZ=V_Z^*V_Z\);
\item a completed carrier feature map
      \[
      W_{\mathrm{cmp}}=
      \begin{bmatrix}
      W_{\mathrm{pr}}\\ W_\infty\\ W_{\mathrm{pole}}\\ W_{\mathrm{tail}}
      \end{bmatrix};
      \]
\item equality of the Weil quadratic form with
      \[
      y\mapsto \langle y,(\Kcmp-\AZ)y\rangle
      \]
      on a dense or exhaustive test class;
\item closure of this positivity from the test class to the completed response space;
\item a positive feature decomposition of the signed prime/gamma/pole/tail terms, or an accepted defect record
      \[
      \AZ\preceq\Kcmp+E_{\mathrm{EF}};
      \]
\item a no-smuggling record: \(W_{\mathrm{cmp}}\) must be built from upstream carrier data, not from target zero displacement.
\end{enumerate}
\end{obligation}

\section{Why scalar or partial positivity is insufficient}

\begin{warning}[Trace equality]
Trace equality
\[
        \operatorname{tr}\AZ=\operatorname{tr}\Kcmp
\]
does not imply \(\AZ\preceq\Kcmp\).
\end{warning}

For example, on \(\mathbb R^2\),
\[
        \AZ=\begin{pmatrix}2&0\\0&0\end{pmatrix},
        \qquad
        \Kcmp=\begin{pmatrix}1&0\\0&1\end{pmatrix}
\]
have the same trace but \(\AZ\npreceq\Kcmp\).

\begin{warning}[Cone or subspace-only testing]
If positivity is checked only on a proper cone or a proper subspace, it need not
imply \(\AZ\preceq\Kcmp\) on the full anti-invariant response space.
\end{warning}

For example
\[
        D=\Kcmp-\AZ=
        \begin{pmatrix}1&0\\0&-1\end{pmatrix}
\]
is positive on the one-dimensional test subspace spanned by \(e_1\), but is not
positive on the full space.

\begin{warning}[Signed components]
A signed explicit formula is not yet a feature-map certificate.  The framework
needs positive carrier-side maps \(W_{\mathrm{pr}},W_\infty,W_{\mathrm{pole}},
W_{\mathrm{tail}}\), or an accepted completed grouping whose total is positive.
\end{warning}

\section{Relation to classical Weil positivity}

The classical Weil criterion may be read as saying that a certain completed
explicit-formula quadratic form is positive for all admissible tests if and only
if RH holds.  In the framework, this becomes:

\[
        Q_W(y)=\langle y,(\Kcmp-\AZ)y\rangle.
\]

If this identity is established on the exact anti-invariant response space and
\(Q_W\ge0\) there, then the Douglas gate follows.  Conversely, the Douglas gate
is a constructive operator-valued strengthening/packaging of the same
positivity: it not only says the residual is positive, but asks for the carrier
side to be expressed by declared positive features and for zero displacement to
factor through those features without amplification.

\section{Three statuses}

\begin{description}[leftmargin=2.8cm]
\item[\texttt{Weil-positive-full}] \(Q_W=\langle\cdot,(\Kcmp-\AZ)\cdot\rangle\) on the full completed response space and \(Q_W\ge0\).  This implies Douglas domination.
\item[\texttt{Weil-positive-test}] Positivity is known only on an admissible test class.  This is support unless density/closure and tail records promote it.
\item[\texttt{signed-EF-only}] A signed explicit formula has been written, but no positive feature decomposition or domination has been supplied.  This is not an O1 bridge.
\end{description}

\section{Conclusion}

Classical Weil positivity is not irrelevant to the Douglas gate; it is the
closest known analytic-number-theoretic shape.  But the gate requires a specific
upgrade:
\[
        \text{Weil positivity}
        \quad\leadsto\quad
        \text{operator Loewner domination}
        \quad\leadsto\quad
        \text{Douglas contraction}
        \quad\leadsto\quad
        \text{positive completed feature maps}.
\]
The remaining RH work is to earn this upgrade on the actual completed zeta
carrier, without replacing signed explicit-formula accounting by an unaudited
positive shadow.
\end{document}
